% \documentclass{article}

% \usepackage{xcolor} % Required for the colors

% \begin{document}

% The \dirtree command uses periods followed by a number to indicate depth.
% Always end each line with a period!
\dirtree{%
.1 \textbf{Participación ciudadana: Áreas Temáticas y Competencias}. %[cite: 6]
.2 \textbf{Movilidad y Accesibilidad}. %[cite: 6]
.3 \textcolor{magenta}{Competencia del Cabildo de Tenerife:}. %[cite: 6]
.4 Transporte público nocturno. %[cite: 6]
.4 Transporte público más barato para reducir el tráfico. %[cite: 6]
.4 No a la Vía Exterior. %[cite: 6]
.4 Guagua circular en el Suroeste. %[cite: 6]
.4 Aceras en TF-180 Vistabella. %[cite: 6]
.4 Miradores en Anaga. %[cite: 6]
.4 Mejora de la frecuencia del transporte público. %[cite: 6]
.4 Apartaderos carreteras Anaga. %[cite: 6]
.3 \textcolor{cyan}{Competencia del Ayuntamiento:}. %[cite: 6]
.4 Solución al estrechamiento carretera de acceso al Parque Rural por San Andrés. %[cite: 6]
.4 Conexión viaria entre barrios del Suroeste. %[cite: 6]
.4 Peatonalización de los centros urbanos. %[cite: 6]
.4 Ciudad para el peatón, no para el coche. %[cite: 6]
.4 Restringir aparcamientos en el centro SC (franja horaria). %[cite: 6]
.4 Determinar recorridos preferentes que sean accesibles a sillas de ruedas. %[cite: 6]
.2 \textbf{Espacio Libre Público}. %[cite: 6]
.3 \textcolor{cyan}{Competencia del Ayuntamiento:}. %[cite: 6]
.4 Uso de barrancos como zonas verdes. %[cite: 6]
.4 No a la privatización del espacio público (terrazas). %[cite: 6]
.4 Ejecución de los Espacios Libres previstos. %[cite: 6]
.4 Rutas deportivas iluminadas por la noche. %[cite: 6]
.4 Rutas deportivas segregadas y alejadas del tráfico. %[cite: 6]
.4 Renaturalización espacios urbanos. %[cite: 6]
.4 Apertura al mar en todo el municipio. %[cite: 6]
.2 \textbf{Dotaciones, Equipamientos e Infraestructuras}. %[cite: 6]
.3 \textcolor{cyan}{Competencia del Ayuntamiento:}. %[cite: 6]
.4 Pequeños intercambiadores modales en el perímetro de los barrios para dejar coche. %[cite: 6]
.4 Centros Dotacionales Integrales polivalentes. %[cite: 6]
.4 Centro de Salud en el Suroeste. %[cite: 6]
.4 Centros de participación cívica. %[cite: 6]
.4 {Ejecución dotaciones previstas en Acorán. Evitar ciudad dormitorio}. %[cite: 6]
.4 Mercadillo del Agricultor Suroeste. %[cite: 6]
.4 Camping entorno Las Teresitas. %[cite: 6]
.2 \textbf{Patrimonio Cultural}. %[cite: 6]
.3 \textcolor{cyan}{Competencia del Ayuntamiento:}. %[cite: 6]
.4 Protección del patrimonio inmaterial. %[cite: 6]
.4 Preservación de los estanques de agua para las aves migratorias y la agricultura. %[cite: 6]
.4 Revisar el Catálogo de Protección para ajustarlo a la realidad. %[cite: 6]
.4 Protección de eras, bancales, paseros. %[cite: 6]
.4 Rehabilitación del Barranco del Muerto. %[cite: 6]
.3 \textcolor{green}{Plan Estratégico Municipal:}. %[cite: 6]
.4 Fomentar comercio de proximidad. %[cite: 6]
.4 Ayudas a la Rehabilitación Edificios. %[cite: 6]
.4 Apostar por un parque de vivienda de alquiler. %[cite: 6]
.4 Divulgación del patrimonio material e inmaterial. %[cite: 6]
.2 \textbf{Desarrollo Sostenible}. %[cite: 6]
.3 \textcolor{cyan}{Competencia del Ayuntamiento:}. %[cite: 6]
.4 Sistemas bio urbanos, calles orientadas a corrientes de aire para evitar aires acondicionados. %[cite: 6]
.4 Rehabilitar el paisaje de la Montaña Taco. %[cite: 6]
.4 Creación de un parque agroecológico en el Suroeste. %[cite: 6]
.4 Huertos urbanos. %[cite: 6]
.4 Separar industrias molestas del área residencial. %[cite: 6]
.3 \textcolor{green}{Plan Estratégico Municipal:}. %[cite: 6]
.4 Lucha contra Cambio Climático. %[cite: 6]
.4 Fomento de las Energías Limpias. %[cite: 6]
.4 Hacer atractiva la ciudad para residir y no perder población. %[cite: 6]
}

% \end{document}